\section{Retrieval Architectures for Legal Intelligence}
\label{sec:retrieval}

\subsection{Why retrieval is a distinct legal AI problem}

In a legal RAG system, the generator cannot reliably compensate for evidence that the retrieval
component fails to identify. If the relevant statutory provision, exception, or judicial
interpretation is absent from the retrieved context, the model may generate an incomplete answer even
when its language-generation capabilities are strong.

Legal retrieval therefore requires more than finding documents that discuss the same topic. The
system must identify evidence that is relevant to the user's legal question, sufficiently
authoritative for the intended purpose, and detailed enough to support the answer. Depending on the
query, this may require an exact statutory citation, a particular paragraph in a judgment, a related
provision, or several authorities that must be considered together.

Research on legal case retrieval has highlighted the specialized nature of this task and the need to
evaluate retrieval methods using legal-specific datasets and criteria \citep{feng2024}. Four
dimensions are especially important:
\begin{itemize}[leftmargin=1.4em, itemsep=2pt, topsep=2pt]
\item \textbf{Lexical relevance:} Does the evidence contain the exact legal terms, names, identifiers,
or expressions needed?
\item \textbf{Semantic relevance:} Does the evidence address the underlying legal issue, even if it
uses different wording?
\item \textbf{Structural relevance:} Does the retrieved passage retain the context of the provision,
document section, or judgment from which it originates?
\item \textbf{Evidential sufficiency:} Does the retrieved material adequately support the intended
answer, including relevant exceptions or qualifications?
\end{itemize}

These dimensions are related but not interchangeable. A passage may be semantically relevant without
containing the exact provision required. An exact match may retrieve a provision without its
exception. A correctly retrieved paragraph may still be insufficient to establish the legal
conclusion.

\subsection{Lexical retrieval: precision through exact wording}

Lexical retrieval ranks documents or passages according to their textual relationship with the query.
One widely used method is BM25, a probabilistic ranking function that considers term frequency,
inverse document frequency, and document length.

In simplified terms, BM25 gives greater importance to query terms that are informative across the
collection and evaluates how frequently those terms occur in a candidate document, with adjustments
that prevent very long documents from receiving an unfair advantage merely because they contain more
words.

This is useful in legal retrieval because legal questions frequently include highly specific
expressions. A search for a section number, statutory phrase, party name, case citation, or defined
legal term may require exact matching that a purely semantic search does not consistently prioritize.

For example, a query that refers to a particular article of a constitution may benefit from matching
the article identifier directly. A dense embedding system might rank a conceptually similar passage
above the exact provision if the two passages are semantically close.

Lexical retrieval also has limitations. Users may describe a legal concept in everyday language while
the relevant authority uses formal terminology. A query about a right to a fair hearing, for instance,
may use different wording from the authoritative provision or the language of a judgment. A lexical
system can miss relevant evidence when the query and the source use different expressions.

BM25 should therefore be understood as a strong retrieval baseline, particularly for exact legal
terminology, rather than a complete solution to legal information retrieval.

\subsection{Dense retrieval: matching legal meaning}

Dense retrieval represents queries and documents as numerical vectors, usually produced by an
embedding model. The system retrieves passages whose vectors are close to the query vector under a
chosen similarity or distance measure.

The advantage is that retrieval can respond to conceptual similarity even when the wording differs. A
user does not necessarily need to know the exact phrase used in a statute or judgment to retrieve
potentially relevant material. This is particularly useful for natural-language questions, paraphrases,
and legal concepts expressed differently across documents.

However, semantic similarity is not equivalent to legal relevance. Two provisions may discuss the same
broad topic but apply under different conditions. A judgment may mention a legal doctrine while
rejecting its applicability to the facts. A dense retriever may rank such material highly because the
language is similar, even though the material does not support the answer.

Dense retrieval also depends on the quality of the embedding model and its training data. A model may
represent general semantic relationships well while failing to capture distinctions involving legal
terminology, jurisdiction, procedural posture, or the authority of a source.

The practical implication is that dense retrieval is valuable for discovering conceptually related
evidence, but its results must still be assessed for legal relevance and sufficiency.

\subsection{Hybrid retrieval: combining lexical and semantic evidence}

Lexical and dense retrieval address complementary weaknesses. Lexical search is often effective for
exact terms, identifiers, and citations; dense search can retrieve conceptually relevant passages
expressed using different language. Hybrid retrieval combines these signals to reduce dependence on
either one alone.

A hybrid architecture may run BM25 and dense retrieval in parallel, obtain ranked lists from both, and
combine them using a fusion method. One widely used method is Reciprocal Rank Fusion (RRF), introduced
by \citet{cormack2009}. RRF combines rankings using the reciprocal of each document's rank, with a
constant controlling how much the rank position contributes to the final score.

Its central advantage is that it can combine rankings without requiring their original scores to be
directly comparable. BM25 scores and dense similarity scores may have different distributions; rank
fusion provides a way to combine their relative ordering instead. Conceptually, RRF can be written as
\begin{equation}
\operatorname{RRF}(d)=\sum_{r \in R}\frac{1}{k+\operatorname{rank}_{r}(d)},
\end{equation}
where $d$ is a candidate document or passage, $R$ is the set of ranked retrieval results,
$\operatorname{rank}_{r}(d)$ is the rank of the candidate in retrieval list $r$, and $k$ is a
smoothing constant.

A passage ranked highly by multiple retrieval methods can receive a strong fused score. A passage
found by only one method may still rank well if that method places it sufficiently high.

Importantly, RRF combines rankings; it does not independently verify legal correctness. If both
retrieval streams return irrelevant evidence, fusion cannot turn that evidence into authoritative
support. Likewise, RRF does not automatically preserve statutory hierarchy or determine whether a
retrieved passage contains a decisive exception. Table~\ref{tab:retrstrat} compares the principal
retrieval strategies.

\begin{table*}[tp]
\centering
\caption{Comparison of principal retrieval strategies.}
\label{tab:retrstrat}
\footnotesize
\setlength{\tabcolsep}{5pt}
\renewcommand{\arraystretch}{1.3}
\begin{tabularx}{\textwidth}{>{\raggedright\arraybackslash}p{2.9cm} >{\raggedright\arraybackslash}X >{\raggedright\arraybackslash}X >{\raggedright\arraybackslash}X}
\toprule
\textbf{Strategy} & \textbf{Primary strength} & \textbf{Main weakness} & \textbf{Role in legal RAG} \\
\midrule
BM25 / lexical retrieval & Exact legal terms, names, identifiers, and statutory references. & May miss relevant passages expressed using different terminology. & Strong baseline for exact-match legal search. \\
Dense retrieval & Semantic similarity and paraphrase matching. & Can retrieve conceptually similar but legally inapplicable passages. & Expands retrieval beyond exact wording. \\
Hybrid retrieval & Combines exact-term matching and semantic search. & Requires tuning and does not guarantee evidence correctness. & Broadens the candidate evidence set. \\
Reciprocal Rank Fusion & Combines ranked lists without directly comparing their raw scores. & Cannot correct poor candidate lists or assess legal authority by itself. & Integrates complementary retrieval results. \\
Reranking & Reassesses candidate passages using a more targeted relevance model. & Adds computation and remains dependent on the quality of candidates and training. & Refines the ordering before evidence is passed to the generator. \\
\bottomrule
\end{tabularx}
\end{table*}

Hybrid retrieval is therefore best understood as a method for improving evidence discovery. Whether it
improves final answer reliability must be tested separately through retrieval metrics and
grounded-generation evaluation.

\subsection{Reranking and evidence selection}

Retrieval systems often return a relatively broad set of candidate passages. These candidates may be
sufficient for an initial search but unsuitable for direct inclusion in a generator's context.
Reranking introduces a second-stage relevance assessment that reorders candidates according to a more
targeted relationship between the query and each passage.

A typical architecture first uses a relatively efficient retriever to obtain candidate passages. A
reranker then evaluates the query--passage pairs more carefully and assigns a refined ranking. Only
the highest-ranked evidence is passed to the generator.

In legal applications, reranking can help distinguish between a passage that merely mentions a legal
concept and one that directly addresses the question. Nevertheless, relevance alone is not enough. A
reranker trained to estimate textual relevance may still fail to recognize that a provision has been
superseded, that a judgment is distinguishable, or that a qualification elsewhere in the statute
changes the interpretation.

Evidence selection must therefore consider more than similarity. Where available, useful signals
include source authority, jurisdiction, date or version, document type, and structural relationships.
These signals can be incorporated into filtering, ranking, or subsequent verification, but their
effectiveness depends on the quality of the underlying metadata and rules.

A further issue is context-window management. Providing every retrieved passage to the generator can
introduce noise, redundancy, or conflicting information. Providing too little context can omit an
exception or related provision. The goal is not simply to maximize the number of retrieved passages;
it is to provide a compact and sufficiently complete set of evidence for the question.

\subsection{Structure-aware indexing: preserving legal hierarchy}

Legal documents are inherently hierarchical. A statute may contain parts, chapters, sections,
subsections, clauses, and provisos. A judgment may contain procedural history, facts, issues,
arguments, reasoning, and the final decision. These structures help determine how an individual
passage should be interpreted.

Conventional chunking strategies often divide documents according to a fixed token limit or a simple
text boundary. Such methods are computationally convenient, but they can separate a provision from its
heading, remove the parent section from a subsection, or split an exception away from the rule it
qualifies. Consider a simplified example:

\begin{quote}
\textit{Illustrative legal hierarchy (schematic):}
\begin{itemize}[leftmargin=1.3em, itemsep=1pt, topsep=2pt]
\item \textbf{Act}: the governing legislative document
  \begin{itemize}[leftmargin=1.2em, itemsep=1pt, topsep=1pt]
  \item \textbf{Section 10}: general legal rule
    \begin{itemize}[leftmargin=1.2em, itemsep=1pt, topsep=1pt]
    \item \textbf{Section 10(2)}: condition or qualification
    \item \textbf{Section 10(3)}: exception or additional requirement
    \end{itemize}
  \end{itemize}
\end{itemize}
The example is schematic, not a statement about a particular statute.
\end{quote}

If a system retrieves only the general rule, it may overlook a condition or exception in a subordinate
provision. The retrieved text may be relevant yet insufficient to support a complete answer.
Structure-aware indexing seeks to reduce this risk by preserving information about the location and
relationships of passages within the original document.

A structure-aware index may retain hierarchical breadcrumbs, section identifiers, headings,
parent-child relationships, document metadata, and source references alongside the text itself. For
example, a passage can be associated with a breadcrumb such as Constitution $\to$ Part~II $\to$
Chapter~1 $\to$ Article~10A. When retrieved, this information helps the system identify where the
passage belongs and may support retrieval of its surrounding context.

Structure-aware indexing is not identical to semantic chunking. Semantic chunking attempts to create
text segments that are coherent in meaning; structure-aware indexing additionally represents the
passage's position within the source document. The two can be combined: a system may create
semantically coherent chunks while retaining their statutory hierarchy and provenance.

Nor does preserving hierarchy automatically resolve legal interpretation. The system must still
determine which related provisions matter to the question, whether the source is authoritative and
current, and whether the combined evidence supports the conclusion. Structure-aware indexing supplies
useful context; it does not replace legal reasoning or verification.

\subsection{Knowledge-graph-based retrieval}

Knowledge graphs offer another way to represent relationships among legal entities and authorities.
Rather than storing documents only as independent passages, a graph represents entities as nodes and
relationships as edges. Legal entities may include statutes, provisions, cases, courts, parties, legal
concepts, and amendments. A graph can represent relationships such as:
\begin{itemize}[leftmargin=1.4em, itemsep=1pt, topsep=2pt]
\item a provision belonging to a particular statute;
\item a judgment interpreting or citing a provision;
\item an amendment modifying an earlier provision;
\item a court issuing a decision;
\item a judgment addressing a defined legal issue.
\end{itemize}

Such relationships can support retrieval beyond direct text similarity. A system may retrieve a
provision and then traverse relevant connections to locate judgments that interpret it or related
provisions that constrain its application. When combined with dense or lexical retrieval, graph
traversal can help connect evidence that is dispersed across documents.

However, graph quality depends on the accuracy of entity extraction, relationship identification,
source linking, and temporal information. An incorrectly extracted citation can create a misleading
connection. A graph that records that a judgment cites a provision does not necessarily establish that
the judgment adopts, applies, or endorses the provision's interpretation.

Knowledge graphs and structure-aware indexes also solve different problems. A hierarchical index
preserves a document's internal organization; a knowledge graph can represent relationships across
multiple documents. They can be complementary components of a retrieval system rather than competing
alternatives.

\subsection{A unified retrieval architecture for legal intelligence}

The preceding approaches address different parts of the retrieval problem. Lexical retrieval identifies
exact textual matches; dense retrieval expands discovery to semantically related passages; rank fusion
combines their results; reranking refines candidate ordering; structure-aware indexing preserves
document context; and knowledge graphs represent explicit relationships among legal authorities. A
unified architecture can combine these components in a staged pipeline, illustrated in
Figure~\ref{fig:arch}.

% Conceptual unified legal retrieval architecture (vertical staged pipeline). Vector TikZ.
\begin{figure}[t]
\centering
\resizebox{\columnwidth}{!}{%
\begin{tikzpicture}[
  font=\footnotesize,
  stage/.style={draw, thick, rounded corners=2pt, align=center, text width=6.6cm,
                minimum height=8mm, inner sep=3pt, fill=legalblue!6, draw=legalblue!70},
  sub/.style={draw, thick, rounded corners=2pt, align=center, text width=3.1cm,
              minimum height=8mm, inner sep=3pt, fill=black!4, draw=black!55},
  fl/.style={-{Latex[length=2.4mm]}, draw=black!50, line width=0.8pt}]

  \node[stage] (ing) at (0,0) {\textbf{1. Legal document ingestion}\\\scriptsize statutes, regulations, judgments, and metadata};
  \node[stage] (pre) at (0,-1.5) {\textbf{2. Structure-aware preprocessing}\\\scriptsize hierarchical segmentation, metadata, source links, provenance};
  \node[sub] (lex) at (-1.85,-3.1) {\textbf{Lexical retrieval}\\\scriptsize BM25, exact legal expressions};
  \node[sub] (den) at ( 1.85,-3.1) {\textbf{Dense retrieval}\\\scriptsize embedding-based semantic similarity};
  \node[stage] (fus) at (0,-4.7) {\textbf{3. Rank fusion and reranking}\\\scriptsize combine candidate lists and refine relevance};
  \node[stage] (ctx) at (0,-6.2) {\textbf{4. Context and authority checks}\\\scriptsize parent provisions, related authorities, jurisdiction, provenance};
  \node[stage] (gen) at (0,-7.7) {\textbf{5. Evidence-grounded generation and verification}\\\scriptsize generate an answer and check its claims against the sources};

  \draw[fl] (ing) -- (pre);
  \draw[fl] (pre.south) -- (lex.north);
  \draw[fl] (pre.south) -- (den.north);
  \draw[fl] (lex.south) -- (fus.north);
  \draw[fl] (den.south) -- (fus.north);
  \draw[fl] (fus) -- (ctx);
  \draw[fl] (ctx) -- (gen);
\end{tikzpicture}}
\caption{Conceptual architecture synthesizing complementary legal retrieval methods. This is a
proposed reference architecture, not a claim that every component has been validated together in a
single experiment.}
\label{fig:arch}
\end{figure}


A knowledge graph can be connected to the context-and-authority stage to discover relationships among
retrieved sources. An agentic controller may also invoke retrieval again when the evidence is
incomplete or contradictory. Both additions should be justified by measurable benefits rather than
included merely to increase architectural complexity.

\subsection{How retrieval architectures should be evaluated}

A retrieval architecture should be assessed independently from the quality of its generated answers.
If the retrieval component does not return the required evidence, the generator's output cannot
reliably compensate for that failure. Conversely, successful retrieval does not establish that the
generator will interpret the evidence correctly. Several complementary metrics are useful, as
summarized in Table~\ref{tab:metrics}.

\begin{table*}[tp]
\centering
\caption{Complementary metrics for evaluating legal retrieval architectures.}
\label{tab:metrics}
\footnotesize
\setlength{\tabcolsep}{5pt}
\renewcommand{\arraystretch}{1.3}
\begin{tabularx}{\textwidth}{>{\raggedright\arraybackslash}p{4.2cm} >{\raggedright\arraybackslash}X >{\raggedright\arraybackslash}X}
\toprule
\textbf{Metric} & \textbf{What it measures} & \textbf{Why it matters in legal retrieval} \\
\midrule
Recall@$k$ & Whether relevant evidence appears among the top $k$ retrieved results. & Measures the system's ability to find needed authorities or passages. \\
Precision@$k$ & The proportion of the top $k$ results that are relevant. & Indicates how much irrelevant material may enter the generation context. \\
Mean Reciprocal Rank (MRR) & The reciprocal rank of the first relevant result, averaged across queries. & Measures how early the system finds a relevant passage. \\
Normalized Discounted Cumulative Gain (nDCG) & Ranking quality when results have graded relevance. & Helps distinguish highly relevant legal evidence from merely related passages. \\
Citation precision & The proportion of supplied citations that correctly support their associated claims, under a defined annotation protocol. & Measures whether references are actually evidentially useful. \\
Citation recall & The proportion of required supporting sources or claims that receive appropriate citations, under a defined protocol. & Helps detect missing legal support. \\
Unsupported-claim rate & The proportion of evaluated claims that lack adequate support in the permitted evidence. & Evaluates grounding beyond the presence of citations. \\
Latency and cost & Time and computational or service cost per query. & Helps determine whether architectural improvements are practical. \\
\bottomrule
\end{tabularx}
\end{table*}

These metrics require clear annotation guidelines. For example, relevance may be graded at the document
level or the passage level; the two can produce different results. Similarly, citation correctness
requires an explicit definition of what counts as support. A citation that mentions the same topic is
not necessarily a citation that substantiates the claim.

A rigorous evaluation should therefore use a consistent query set and legal corpus, compare lexical,
dense, hybrid, and structure-aware configurations under controlled conditions, and measure both
retrieval and final answer reliability. Where feasible, ablation studies should remove one component at
a time to determine whether improvements can be attributed to rank fusion, structural metadata,
reranking, graph traversal, or another architectural choice.

\subsection{Synthesis and remaining research challenges}

The development of legal retrieval architectures reflects a progression from finding text that matches
a query toward identifying evidence that is relevant, structurally contextualized, and suitable for
supporting a legal conclusion. Lexical retrieval provides a strong foundation for exact legal
terminology. Dense retrieval addresses differences in wording. Hybrid retrieval and rank fusion combine
complementary signals, while reranking refines the candidate set. Structure-aware indexing preserves
the organization of legal documents, and knowledge graphs can connect authorities across documents.

None of these components independently guarantees legal correctness. Their value depends on the legal
corpus, source provenance, retrieval configuration, quality of metadata, and evaluation protocol. A
sophisticated retriever can still return an incomplete set of authorities, while a well-ranked passage
can remain insufficient if its parent provision or relevant exception is missing.

The resulting research direction is not simply to build a larger or more complicated retrieval
pipeline. It is to design systems in which the evidence required for a legal answer can be retrieved,
contextualized, and verified systematically. This provides the technical foundation for the next part
of the survey, which examines hallucination, factual support, and citation correctness as separate but
connected dimensions of legal AI reliability.
